% Transcribed from CMM233 Summary.docx.
\IfFileExists{main.tex}{\documentclass[main.tex]{subfiles}}{
  \IfFileExists{../../main.tex}{\documentclass[../../main.tex]{subfiles}}{
    \IfFileExists{../main.tex}{\documentclass[../main.tex]{subfiles}}{
      \documentclass[Book\space NL/main.tex]{subfiles}
    }
  }
}
\begin{document}
\ifSubfilesClassLoaded{
  \ifcsname appendixchapterparent\endcsname
  \else
    \standaloneappendixchaptersetup{D}{1}
  \fi
}{}
\edef\reviewnotesimagepath{\booknlroot appendices/appendix 4-Review Notes/images/}
\chapter{Cardiovascular and Respiratory System Review}
\labch{appendix-d1-cardiovascular-respiratory-system-review}
\margintoc

\begin{center}\Large\textbf{Systems Summary - Tai’s edition}\end{center}

\section{Cardiovascular Anatomy and Physiology}
\subsection{Thorax, Heart, and Mediastinum}
\begin{itemize}
\item Compartments of the thoracic cavity:
\begin{itemize}
\item Pleural cavities (left and right) -- cavities where the lungs are located
\item Mediastinum -- the central compartment containing the heart
\end{itemize}
\item Thoracic wall (bony thorax) -- formed by the sternum, ribs, and thoracic vertebrae
\begin{itemize}
\item The sternum consists of, from top to bottom, the manubrium, body, and xiphoid process.
\item There are 12 pairs of ribs: ribs 1--7 are true ribs, and ribs 8--12 are false ribs.
\item[]
\begin{itemize}
\item True ribs articulate with their own costal cartilages and the sternum.
\item False ribs articulate with the costal cartilage of the rib above (ribs 8--10) or are floating ribs (ribs 11--12).
\item Ribs 1--2 and 11--12 are atypical in morphology.
\end{itemize}
\item Typical rib structures, from posterior to anterior:
\item[]
\begin{itemize}
\item Head with superior and inferior articular facets.
\item Neck.
\item Tubercle, with an articular facet for the transverse process of a vertebra.
\item Angle.
\item Costal groove.
\end{itemize}
\end{itemize}
\end{itemize}
\noindent The typical rib and its major landmarks are illustrated in Figure~\ref{fig:appendix-figure-d1-thoracic-anatomy-typical-rib-landmarks}.
\begin{center}\includegraphics[width=0.55\textwidth,height=0.42\textheight,keepaspectratio]{\reviewnotesimagepath d1-figure-thoracic-anatomy-typical-rib-landmarks.png}\par\captionof{figure}{Posterior-to-anterior landmarks on a typical rib, including its articular facets, tubercle, angle, and costal groove.}\labfig{appendix-figure-d1-thoracic-anatomy-typical-rib-landmarks}\end{center}
\begin{itemize}
\item[]
\begin{itemize}
\item[]
\begin{itemize}
\item Atypical rib structures.
\begin{itemize}
\item Rib I is flat and broad, with superior and inferior surfaces, a single articular facet, a scalene tubercle, and grooves for the subclavian vessels.
\item Rib II is similar to rib I, but is longer and has two articular facets.
\item Ribs XI and XII are short, have little curvature, and are pointed anteriorly. They lack a neck and tubercle; rib XI has a shallow costal groove, whereas rib XII has none.
\end{itemize}
\end{itemize}
\end{itemize}
\end{itemize}
\noindent The atypical forms of ribs I, II, XI, and XII are compared in Figures~\ref{fig:appendix-figure-d1-thoracic-anatomy-atypical-rib-comparison} and~\ref{fig:appendix-figure-d1-thoracic-anatomy-floating-rib-landmarks}.
\begin{center}\includegraphics[width=0.58\textwidth,height=0.46\textheight,keepaspectratio]{\reviewnotesimagepath d1-figure-thoracic-anatomy-atypical-rib-comparison.png}\par\captionof{figure}{The first and second ribs differ from a typical rib in their broad superior surfaces and distinctive landmarks.}\labfig{appendix-figure-d1-thoracic-anatomy-atypical-rib-comparison}\end{center}
\begin{center}\includegraphics[width=0.72\textwidth,height=0.36\textheight,keepaspectratio]{\reviewnotesimagepath d1-figure-thoracic-anatomy-floating-rib-landmarks.png}\par\captionof{figure}{The eleventh and twelfth ribs are short, minimally curved, and lack the neck and tubercle found on typical ribs.}\labfig{appendix-figure-d1-thoracic-anatomy-floating-rib-landmarks}\end{center}
\begin{itemize}
\item[]
\begin{itemize}
\item Structure of thoracic vertebrae:
\item[]
\begin{itemize}
\item Vertebrae T1--T12.
\item Heart-shaped vertebral body with superior and inferior costal facets for articulation with the ribs.
\item Transverse process that projects posterolaterally.
\item Circular vertebral foramen.
\item Long spinous process projecting posteriorly.
\end{itemize}
\item Muscles of the thorax:
\item[]
\begin{itemize}
\item External intercostals originate from the lower margin of the rib above and insert on the superior margin of the rib below. They raise the ribs during inspiration.
\item Internal intercostals run between adjacent ribs. Their interosseous portion generally depresses the ribs during forced expiration, while the interchondral portion can elevate them during inspiration.
\item Innermost intercostals originate from the subcostal groove of the rib above and insert on the upper border of the rib below. They lower the ribs during expiration.
\item All intercostals are innervated by their respective intercostal nerves (T1--T11). They support the intercostal spaces and stabilize the chest wall.
\item Subcostal muscles depress the ribs. They originate from the lower margins of lower ribs and insert on the inner surfaces of ribs two or three levels below. They are innervated by various intercostal nerves.
\item Transversus thoracis arises from the posterior sternum and xiphoid process and inserts on the internal surfaces of costal cartilages, usually ribs 2--6. It weakly depresses the ribs and is innervated by intercostal nerves.
\end{itemize}
\item Contents of the intercostal spaces: intercostal muscles, intercostal nerves, and intercostal vessels (veins and arteries).
\end{itemize}
\item Vessels and innervation of the intercostals:
\begin{itemize}
\item Subclavian artery \(\rightarrow\) internal thoracic artery \(\rightarrow\) anterior intercostal arteries (upper six ribs).
\item Internal thoracic artery \(\rightarrow\) musculophrenic artery \(\rightarrow\) anterior intercostal arteries (lower six ribs).
\item Thoracic aorta \(\rightarrow\) posterior intercostal arteries.
\item Anterior intercostal veins \(\rightarrow\) internal thoracic veins \(\rightarrow\) brachiocephalic veins \(\rightarrow\) superior vena cava.
\item Posterior intercostal veins \(\rightarrow\) azygos vein \(\rightarrow\) superior vena cava.
\item Intercostal nerves are the anterior rami of thoracic spinal nerves T1--T11. They give off lateral cutaneous branches and terminate as anterior cutaneous branches; these cutaneous branches divide into anterior and posterior or medial and lateral branches, respectively.
\item T12 arises as the subcostal nerve and has motor and sensory functions.
\end{itemize}
\item Mediastinum.
\end{itemize}
\subsection{Histology of the Cardiovascular System}
\begin{itemize}
\item The circulatory system is divided into cardiovascular and lymphatic systems.
\item The heart pumps blood through pulmonary and systemic circuits: arteries carry blood away from the heart, capillaries are sites of nutrient, gas, and hormone/ligand exchange, and veins return blood to the heart.
\item The heart consists of 3 layers:
\begin{itemize}
\item Endocardium
\item[]
\begin{itemize}
\item Inner layer -- endothelium (simple squamous epithelium) and subendothelium (fibroelastic connective tissue).
\item Middle layer -- connective tissue and smooth muscle.
\item Deep layer -- subendocardium (loose connective tissue) with small blood vessels, nerves, and Purkinje fibers.
\end{itemize}
\item Myocardium -- contains cardiac muscle arranged in complex spirals around chamber orifices; it is thicker in the ventricles.
\item Epicardium
\item[]
\begin{itemize}
\item Visceral layer of the serous pericardium.
\item Composed of mesothelium (simple squamous epithelium) and a subepicardial layer of loose connective tissue containing coronary vessels, nerves, and ganglia.
\item Site of fat storage.
\end{itemize}
\end{itemize}
\end{itemize}
\noindent The layered arrangement described here is shown in Figures~\ref{fig:appendix-figure-d1-cardiovascular-histology-heart-wall-layers} and~\ref{fig:appendix-figure-d1-cardiovascular-histology-endocardium-purkinje-fibers}.
\begin{center}\includegraphics[width=0.62\textwidth,height=0.52\textheight,keepaspectratio]{\reviewnotesimagepath d1-figure-cardiovascular-histology-heart-wall-layers.png}\par\captionof{figure}{A ventricular section relates the endocardial lining to papillary muscle, myocardium, and the epicardial surface with its vessels and adipose tissue.}\labfig{appendix-figure-d1-cardiovascular-histology-heart-wall-layers}\end{center}
\begin{center}\includegraphics[width=0.78\textwidth,height=0.48\textheight,keepaspectratio]{\reviewnotesimagepath d1-figure-cardiovascular-histology-endocardium-purkinje-fibers.png}\par\captionof{figure}{The illustrated tissue field distinguishes the endocardial layers from myocardium and locates Purkinje fibers among supporting tissue and small vessels.}\labfig{appendix-figure-d1-cardiovascular-histology-endocardium-purkinje-fibers}\end{center}
\subsection{Topics Listed Without Notes in the Source}
The supplied summary names the following topics but contains no accompanying notes or figures for them. No additional material has been inferred.
\begin{itemize}
\item Development of the heart and blood vessels.
\item The heart as a pump.
\item Hemostasis and blood coagulation.
\item Hemodynamics and the lymphatic system.
\item Physical examination and stimulation (the source says this material was moved to SUM2, but provides no corresponding notes).
\item Regulation of the cardiovascular system.
\item Heart pathology and vascular pathology.
\item Clinical pharmacology of the cardiovascular system.
\item Hypertension and vascular disease.
\item Cardiovascular infections.
\item Heart failure and shock.
\item Congenital cardiovascular abnormalities.
\item Pathology of the heart and blood vessels.
\end{itemize}
\subsection{Electrical Activity of the Heart and ECG}
\begin{itemize}
\item Heart is controlled by electrical activity \(\rightarrow\) changes ion concentration between ECF and ICF (through ion channels)
\item Majority of water in body = intracellular fluid, followed by interstitial fluid (between cells)
\end{itemize}
The source summarizes typical extracellular and intracellular fluid composition in Table~\ref{tab:appendix-table-d1-cardiovascular-histology-fluid-ion-concentrations}. The image table is transcribed below; values and units are retained as shown.
\begin{center}
\captionof{table}{Typical ionic and selected solute concentrations in extracellular and intracellular fluid, transcribed from the source image.}\labtab{appendix-table-d1-cardiovascular-histology-fluid-ion-concentrations}
\small
\begin{tabular}{p{0.30\textwidth}p{0.28\textwidth}p{0.28\textwidth}}
\hline
\textbf{Substance} & \textbf{Extracellular fluid} & \textbf{Intracellular fluid} \\
\hline
Na\textsuperscript{+} & 142 mEq/L & 10 mEq/L \\
K\textsuperscript{+} & 4 mEq/L & 140 mEq/L \\
Ca\textsuperscript{2+} & 2.4 mEq/L & 0.0001 mEq/L \\
Mg\textsuperscript{2+} & 1.2 mEq/L & 58 mEq/L \\
Cl\textsuperscript{--} & 103 mEq/L & 4 mEq/L \\
HCO\textsubscript{3}\textsuperscript{--} & 28 mEq/L & 10 mEq/L \\
Phosphates & 4 mEq/L & 75 mEq/L \\
SO\textsubscript{4}\textsuperscript{2--} & 1 mEq/L & 2 mEq/L \\
Glucose & 90 mg/dL & 0--20 mg/dL \\
Amino acids & 30 mg/dL & 200 mg/dL (?) \\
Cholesterol, phospholipids, and neutral fat & 0.5 g/dL & 2--95 g/dL \\
\(P_{\mathrm{O}_2}\) & 35 mmHg & 20 mmHg (?) \\
\(P_{\mathrm{CO}_2}\) & 46 mmHg & 50 mmHg (?) \\
pH & 7.4 & 7.0 \\
Proteins & 2 g/dL (5 mEq/L) & 16 g/dL (40 mEq/L) \\
\hline
\end{tabular}
\end{center}
\noindent\textbf{[Verify against the cited source: intracellular amino-acid, P\textsubscript{O\textsubscript{2}}, and P\textsubscript{CO\textsubscript{2}} values are marked uncertain in the original image, and the stated intracellular lipid range may be anomalous.]}
\par\noindent\textit{Source noted in the original image: Guyton, 2016.}
\begin{itemize}
\item Ion channels are found in the plasma membrane and many intracellular organelles. They allow a high rate of transport (10\textsuperscript{6}+ ions per second) down an electrochemical gradient, determined by ion concentration and membrane potential, without requiring ATP.
\item Main ion channels: Na\textsuperscript{+}, K\textsuperscript{+}, and Ca\textsuperscript{2+}.
\end{itemize}
\subsubsection{Cardiac Cells}
\begin{itemize}
\item The heart consists of two main types of cells: cardiomyocytes (mechanical cells in the atria and ventricles) and cardiac pacemaker cells (in the SA and AV nodes).
\item Contractile ventricular myocyte action potentials are described in five numbered phases (0--4):
\begin{itemize}
\item Phase 4 (resting phase): approximately $-90$ mV membrane potential.
\item Phase 0 (depolarization): fast Na\textsuperscript{+} channels open.
\item Phase 1 (initial repolarization): fast sodium channels close and potassium channels open.
\item Phase 2 (plateau phase): calcium and potassium channels open.
\item Phase 3 (rapid repolarization): calcium channels close and fast K\textsuperscript{+} channels open.
\item Return to phase 4.
\end{itemize}
\end{itemize}
\iffalse
\section{Laboratory Review}
\subsection{Laboratory Checklists}
\subsubsection{Cardiovascular Histology}
Checklist:\par
\begin{enumerate}
\item Heart wall
\begin{itemize}
\item Endocardium – Purkinje fibers
\item Myocardium
\item Epicardium
\end{itemize}
\item Large (elastic) artery
\item Medium (muscular) artery
\item Small artery
\item Arteriole
\item Capillary
\begin{itemize}
\item Continuous capillary
\item Sinusoidal capillary
\item Pericyte (optional)
\end{itemize}
\item Venule
\item Small vein
\item Medium vein
\item Large vein
\end{enumerate}
\subsubsection{Respiratory Histology}
*Optional อาจะออกเป็นรูป ไม่มี slide\par
Checklist\par
\begin{itemize}
\item Olfactory epithelium: pseudostratified ciliated columnar (identify cell type)
\begin{itemize}
\item Basal cells (basal lamina)
\item Supporting cells
\item Neuron cell (olfactory receptor cell) \(\rightarrow\) prominent nucleolus
\end{itemize}
\item Respiratory epithelium
\begin{itemize}
\item pseudostratified ciliated columnar with goblet cells \(\rightarrow\) ciliated cuboidal epithelium \(\rightarrow\) simple squamous epithelium. Identify cell type.
\end{itemize}
\item Epiglottis: 2 surfaces
\begin{itemize}
\item Anterior/pharyngeal surface - stratified squamous epithelium
\item Posterior/laryngeal surface - pseudociliated columnar with goblet cells
\end{itemize}
\item Trachea cartilage
\begin{itemize}
\item Hyaline cartilage w/ respiratory epithelium
\end{itemize}
\item Lung
\begin{itemize}
\item Bronchus w/ hyaline cartilage fragment - respiratory epithelium w/ less goblet cells. Bands of smooth muscle
\item Bronchiole - no cartilage, only smooth muscle. Respiratory epithelium (cuboid or ciliated columnar)
\item Alveoli - round shaped sac
\item Cluster of alveoli - alveolar sac
\item Respiratory bronchiole - duct epithelium mix with singular alveoli
\item Alveolar duct
\item Type 1 pneumocyte - squamous, no visible nucleus
\item Type 2 pneumocyte - predominant nucleolus, at corner, cuboidal
\end{itemize}
\end{itemize}
• Nasal mucosa (optional)\par
• Larynx (optional)\par
• Structures of conducting portion:\par
– Trachea, intrapulmonary bronchus\par
– Bronchiole \(\rightarrow\) terminal bronchiole\par
• Structures of respiratory portion:\par
– Respiratory bronchiole\par
– Alveolar duct\par
– Alveolar sac, alveoli\par
• Type I pneumocyte\par
• Type II pneumocyte\par
• Dust cell (optional)\par
\subsection{Gross Anatomy}
\subsubsection{Heart, Great Vessels, and Thoracic Wall}
\textbf{Lab 1}\par
\begin{enumerate}
\item \textbf{Skeleton of the thoracic wall}
\begin{itemize}
\item Sternum
\item Manubrium
\item Sternal angle
\item Xiphoid process
\end{itemize}
\item \textbf{Nerves and blood vessels}
\begin{itemize}
\item Intercostal nerve
\item Internal thoracic artery/vein
\item Thoracoepigastric vein
\item Superior vena cava
\item Inferior vena cava
\item Arch of the aorta
\end{itemize}
\item \textbf{Chambers of the heart}
\begin{itemize}
\item Right atrium and right auricle
\item Right ventricle
\item Left ventricle
\item Left atrium and left auricle
\end{itemize}
\item \textbf{Pericardium}
\begin{itemize}
\item Fibrous pericardium
\item Serous pericardium
\end{itemize}
\item \textbf{Valves}
\begin{itemize}
\item Aorta and aortic valve
\item Pulmonary trunk
\item Pulmonary valve
\end{itemize}
\item \textbf{Coronary arteries}
\begin{itemize}
\item Right coronary artery
\item Left coronary artery
\item Right marginal artery
\item Anterior interventricular artery
\item Posterior interventricular artery
\item Circumflex artery
\end{itemize}
\item \textbf{Cardiac veins and coronary sinus}
\begin{itemize}
\item Coronary sinus
\item Great cardiac vein
\item Middle cardiac vein
\item Small cardiac vein
\item Oblique vein of the atrium
\item Anterior cardiac veins
\end{itemize}
\end{enumerate}
\textbf{Lab 2}\par
\begin{enumerate}
\item \textbf{Thoracic wall}
\item \textbf{Posterior mediastinum}
\begin{itemize}
\item Esophagus
\item Thoracic aorta
\item Azygos vein
\item Hemiazygos vein
\item Thoracic duct
\item Vagus nerve
\item Sympathetic trunk
\item Splanchnic nerve
\end{itemize}
\item \textbf{Pharynx}
\begin{itemize}
\item Nasopharynx
\item Laryngopharynx
\item Oropharynx
\end{itemize}
\item \textbf{Trachea}
\begin{itemize}
\item Cricoid cartilage
\item Carina
\item Right main bronchus
\item Left main bronchus
\end{itemize}
\item \textbf{Pleura}
\begin{itemize}
\item Parietal pleura
\item Visceral pleura
\end{itemize}
\item \textbf{Lungs}
\begin{itemize}
\item Right lung
\item Left lung
\item Oblique fissures
\item Horizontal fissure
\item Cardiac impression
\item Cardiac notch
\item Lingula
\item Pulmonary trunk
\item Pulmonary veins
\item Bronchial arteries
\item Bronchial veins
\end{itemize}
\end{enumerate}
\fi
\end{document}